\title{Direct Preference Optimization: Your Language Model is Secretly a Reward Model}

\begin{abstract}
Large-scale unsupervised language models (LMs) are capable of acquiring extensive world knowledge and exhibit some reasoning abilities, yet their behavior remains challenging to control due to the entirely unsupervised paradigm under which they are trained. To enhance model steerability, current approaches typically gather human-annotated assessments comparing the quality of generated outputs, then adjust the parameters of the unsupervised LM to reflect these human preferences. This is frequently accomplished through reinforcement learning from human feedback (RLHF). Despite its widespread adoption, RLHF involves a multi-stage, often unstable process, beginning with training a reward model to capture human judgments and subsequently updating the LM with reinforcement learning in a way that maximizes the learned reward without straying excessively from its initial parameters. This work proposes a novel reward model parameterization within the RLHF framework, permitting analytic extraction of the corresponding optimal policy. This insight enables a reformulation of the RLHF optimization problem into a straightforward classification task. We introduce \textit{Direct Preference Optimization} (DPO), a stable, efficient, and resource-friendly algorithm that obviates the need for language model sampling during adaptation and significantly reduces dependence on hyperparameter optimization. Empirical results indicate that DPO is at least as effective as prior approaches in aligning LMs to human preferences. In particular, DPO not only outperforms PPO-based RLHF in modulating the sentiment of generated text, but also achieves comparable or superior performance in summarization and single-turn dialogue benchmarks, all while maintaining a simpler and more accessible implementation and training procedure.
\end{abstract}

\section{Introduction}
Large-scale unsupervised language models (LMs), when trained on extensive datasets, often develop unexpected proficiencies~\citep{chowdhery2022palm, brown2020language, touvron2023llama,bubeck2023sparks}. Nonetheless, these models are exposed to human-generated data reflecting a diverse spectrum of intentions, priorities, and competencies. Not all of these human characteristics are desirable for the model to emulate; for instance, while an AI-based coding assistant should be capable of \textit{identifying} frequent programming errors for correction purposes, we would ideally prefer the model's code generation to reflect the (possibly infrequent) exemplary programming practices embedded within its training set. In a similar vein, it is valuable for a language model to be \textit{informed} about misconceptions held by a significant portion of the population (e.g., 50\%), yet it is clearly undesirable for the model to repeat or endorse such misconceptions with the same frequency when answering related queries. Thus, it becomes imperative to control the model's \emph{chosen outputs and actions} distinct from its broad \textit{knowledge and competencies}, so as to build AI systems that are reliable, efficient, and amenable to oversight \citep{ouyang2022training}. Traditionally, current approaches rely on reinforcement learning (RL) to align LMs with human preferences; however, we will demonstrate that the optimization goal utilized in these RL-based methods can in fact be achieved directly using a straightforward binary cross-entropy loss, thereby streamlining the preference learning workflow significantly.

Generally, established methods introduce preferred behavioral patterns into language models through carefully selected collections of human preference data that delineate desirable and beneficial behaviors. This stage of learning from preferences follows an initial, extensive unsupervised pre-training phase on large-scale textual corpora. Although a naive method for preference learning involves supervised fine-tuning based on exemplary human-generated responses, reinforcement learning from human or artificial feedback (RLHF/RLAIF; \citep{christiano2017deep,bai2022constitutional}) represents the most prominent family of successful techniques. Within RLHF, a reward model is trained to reflect human preferences, which is then employed in RL to encourage the language model to output responses that receive high reward, all the while preventing significant divergence from the pretrained policy. Despite the substantial advancements RLHF offers in terms of conversation and coding capabilities, this approach remains notably more involved than standard supervised training, necessitating the development and coordination of several language models as well as the inclusion of sampling in the policy optimization loop, which results in considerable computational overhead.

This work presents a method for optimizing language models in alignment with human preferences directly, eliminating the need for explicit reward models or reinforcement learning frameworks. We introduce \textit{Direct Preference Optimization} (DPO), an approach that effectively maximizes the same objective function as traditional RLHF techniques—that is, reward maximization constrained by KL-divergence—while remaining straightforward both in implementation and training procedures. Conceptually, DPO updates the model by amplifying the log-likelihood ratio between preferred and non-preferred outputs, but incorporates adaptive, example-specific importance weights that mitigate model collapse issues observed with naïve probability ratio methods. Like prior approaches, DPO makes use of a theoretical preference framework (e.g., the Bradley-Terry model; \cite{bradley1952rankanalysis}) to evaluate the correspondence between candidate reward functions and observed human preference data. In contrast to conventional methods that leverage the preference model as a loss to train a reward function, subsequently optimizing the policy against this trained reward, DPO reparameterizes the preference loss directly in terms of the policy. Leveraging datasets of human-annotated preferences over model outputs, DPO optimizes the model using a straightforward binary cross entropy loss, ultimately yielding the optimal policy relative to an implicit reward function tailored to the collected preference data.

Our principal contribution lies in introducing Direct Preference Optimization (DPO), a reinforcement learning-free method for adapting language models from preference signals. Empirical results demonstrate that DPO matches or exceeds the performance of existing approaches, including RLHF methods built on PPO, across multiple applications—such as sentiment control, summarization, and conversational modeling—in language models with up to 6B parameters.

Large self-supervised language models have demonstrated the ability to perform some tasks in a zero-shot setting \citep{radford2019language} or by leveraging a few examples in their prompts \citep{gpt3,megatron,chowdhery2022palm}. Nonetheless, fine-tuning these models on datasets containing instructional prompts and human-authored completions has been shown to markedly boost both downstream task effectiveness and user intent alignment \citep{mishra-etal-2022-cross,sanh2022multitask,chung2022scaling,thoppilan2022lamda}. This process, termed `instruction-tuning,' equips large language models (LLMs) with the capability to handle instructions beyond those encountered during fine-tuning, thereby enhancing their general utility \citep{chung2022scaling}. While instruction tuning has been highly successful, collecting \textit{relative} human preference judgments over model responses is generally more scalable than curating expert-level demonstrations. Thus, follow-up research has leveraged human preference datasets to further refine LLMs, leading to advancements in fields such as translation \citep{kreutzer-etal-2018-reliability}, summarization \citep{stiennon2022learning,ziegler2020finetuning}, narrative generation \citep{ziegler2020finetuning}, and task completion \citep{ouyang2022training,ramamurthy2023is}. These approaches typically begin by fitting a neural network reward model to the preference dataset, frequently employing models like Bradley-Terry \citep{bradley1952rankanalysis}, and subsequently use reinforcement learning—such as REINFORCE \citep{williams1992reinforce}, proximal policy optimization (PPO; \cite{schulman2017proximal}), or their modifications \citep{ramamurthy2023is}—to adjust the language model in accordance with the learned reward. Parallel lines of inquiry have employed LLMs tuned with human feedback to produce synthetic preference datasets focusing on targeted features like safety and harmlessness \citep{bai2022constitutional}, using weak forms of supervision via human-authored rubrics to guide LLM annotations. Such methods represent the intersection between research on reinforcement learning-based language model training for diverse targets~\citep{Ranzato2015SequenceLT,paulus2018a,wu2018learning} and frameworks for learning from human preferences \citep{christiano2017deep,kupcsik2018learning}. However, despite the attractiveness of leveraging relative human preferences, the process of reinforcement learning-based fine-tuning for large language models poses considerable practical obstacles; the present work introduces a theoretically-sound method for optimizing relative preferences without reliance on reinforcement learning.

Beyond linguistic contexts, the challenge of inferring policies from preferences has been explored within both bandit and reinforcement learning paradigms, yielding a range of methodological innovations. In the contextual bandit setting, where feedback comes as preferences or action rankings rather than explicit rewards, the task is formalized as the contextual dueling bandit (CDB; \cite{yue2012karmed,dudik2015contextual}). In situations lacking absolute reward signals, theoretical investigations of CDBs replace the optimality criterion with the concept of a \textit{von Neumann winner}: a policy expected to outperform \textit{any} competing policy with probability no less than 50\% \citep{dudik2015contextual}. Notably, CDBs operate with preference labels supplied in an online manner, whereas in the context of learning from human judgments, preference data typically consists of a finite, offline collection of annotated action pairs \citep{yan2022human}. An adjacent line of work, \textit{preference-based RL} (PbRL), utilizes binary preferences that are presumed to arise from some unobserved `scoring' function rather than from known rewards \citep{BusaFekete2014,ruiz2023dueling}. Numerous PbRL algorithms have been developed, with some permitting the reuse of off-policy preference data; these usually proceed by explicitly estimating this hidden scoring function—effectively a reward model—before optimizing policies with respect to it \citep{jain2013learning,BusaFekete2014,christiano2017deep,sadigh2017active,kupcsik2018learning}. In contrast, we introduce a unified policy learning technique that eschews separate reward modeling and instead directly seeks policies that align with expressed preferences.

\section{Preliminaries}\label{section:prelims}

Here we recapitulate the RLHF pipeline as detailed in \citeauthor{ziegler2020finetuning}, and subsequently elaborated by works such as \citep{stiennon2022learning, bai2022training, ouyang2022training}. This workflow can be broken down into three principal stages: (1) supervised fine-tuning (SFT), (2) collecting preference data and fitting a reward model, and (3) policy optimization via reinforcement learning.

\textbf{SFT}: The RLHF process is initiated by adapting a pre-trained language model using supervised learning on task-relevant, high-quality datasets (such as for dialog or summarization), resulting in the model $\pi^\text{SFT}$.

\textbf{Reward Modelling Phase}: Next, the supervised model is supplied with input prompts $x$, generating answer pairs $(y_1, y_2)\sim \pi^\text{SFT}(y \mid x)$. These pairs are evaluated by human annotators, who indicate their preference—denoted $y_w\succ y_l \mid x$—with $y_w$ and $y_l$ referring to the favored and less favored responses from each pair, respectively. These preferences are presumed to reflect an underlying, unobserved reward function $r^*(y, x)$. Various statistical models are available for characterizing such preferences, with the Bradley-Terry (BT) \cite{bradley1952rankanalysis} formulation being widely adopted (and the more general Plackett-Luce ranking models \citep{plackett1975analysis, luce2012individual} applicable where multi-way rankings exist). The BT model assumes that the distribution over human preferences, $p^*$, can be expressed as follows:

Given a static set of pairwise comparisons $\mathcal{D}=\bigl\{x^{(i)}, y_w^{(i)}, y_l^{(i)}\bigr\}_{i=1}^N$ sampled according to $p^*$, we introduce a parameterized reward model $r_{\phi}(x, y)$ and optimize its parameters using maximum likelihood estimation. Recasting the setup as a binary classification problem, the associated loss function is the negative log-likelihood:

\begin{equation}\label{eq:reward_model}
    \mathcal{L}_R(r_{\phi}, \mathcal{D}) = -\mathbb{E}_{(x, y_w, y_l)\sim \mathcal{D}}\bigl[\log \sigma(r_{\phi}(x, y_w)- r_{\phi}(x, y_l))\bigr]
\end{equation}

with $\sigma$ denoting the logistic sigmoid function. For language models, $r_{\phi}(x, y)$ typically uses the SFT model $\pi^\text{SFT}(y \mid x)$ as its base and augments it with a linear head added on top of the last transformer layer, producing a single scalar reward per input-output pair~\cite{ziegler2020finetuning}. To mitigate reward variance, prior works standardize the predicted reward values to ensure $\mathbb{E}_{x,y\sim \mathcal{D}}\left[r_\phi(x, y)\right] = 0$ for all $x$.

\textbf{RL Fine-Tuning Phase}: Once the reward model is trained, it serves as the signal for guiding the language model via reinforcement learning. Building on established methods~\citep{jaques2017sequence, jaques2020human}, the RL objective takes the following form:

\begin{equation}\label{eq:RL}
\max_{\pi_{\theta}}  \mathbb{E}_{x\sim \mathcal{D}, y\sim \pi_{\theta}(y \mid x)}\bigl[r_{\phi}(x, y)\bigr] - \beta\mathbb{D}_{\textrm{KL}}\bigl[\pi_{\theta}(y\mid x)\mid \mid \pi_\text{ref}(y\mid x)\bigr],
\end{equation}

Here, $\beta$ modulates the degree to which the learned policy $\pi_\theta$ diverges from a reference model $\pi_\text{ref}$, typically initialized with the supervised fine-tuned policy $\pi^\text{SFT}$. Initialization of $\pi_\theta$ with $\pi^\text{SFT}$ is common practice. The KL constraint plays a crucial role in confining the model within the domain where the reward estimates are reliable, as well as safeguarding against mode collapse and maintaining output diversity. Since the generative process for sequences is discrete, direct gradient-based optimization is infeasible and RL-based approaches are required. In prevalent approaches~\citep{ziegler2020finetuning, stiennon2022learning, bai2022training, ouyang2022training}, the composite reward is formulated as ${r(x, y) = r_{\phi}(x, y) -\beta (\log \pi_{\theta}(y\mid x) - \log \pi_\text{ref}(y\mid x))}$ and maximized using PPO~\cite{schulman2017proximal}.

\section{Direct Preference Optimization}\label{sec:DPO}

Recognizing the complexity and scale limitations posed by reinforcement learning on large language models, we aim to propose a direct preference-driven alternative for optimizing policies. In contrast to existing RLHF paradigms—which require an explicit reward model to be learned and then separately optimized—our approach capitalizes on a specific structure in the reward parameterization that enables direct, closed-form derivation of the optimal policy. As we explain in subsequent sections, our method hinges on exploiting an analytical correspondence between rewards and their associated optimal policies, thereby re-expressing losses defined over rewards into ones over policies themselves. This change-of-variables perspective eliminates the need for fitting an explicit reward network, while faithfully adhering to established probabilistic preference models, such as the Bradley-Terry formulation. Thus, the policy model concurrently fulfills the roles of language modeling and implicit reward estimation.

\textbf{Derivation of the DPO objective.} We begin with the reinforcement learning (RL) objective defined in Eq.~\ref{eq:RL}, considering an arbitrary reward function $r$, as adopted by previous studies. Building upon earlier work~\citep{peters2007reinforcement, peng2019advantage, korbak2022reinforcement, go2023aligning}, it follows directly that the policy maximizing the KL-regularized reward in Eq.~\ref{eq:RL} assumes the structure:

\begin{equation}\label{eq:op_policy}
    \pi_r(y\mid x) = \frac{1}{Z(x)}\pi_\text{ref}(y\mid x)\exp\left(\frac{1}{\beta}r(x, y)\right),
\end{equation}

where the normalization term, $Z(x) =\sum_{y}\pi_\text{ref}(y\mid x)\exp\left(\frac{1}{\beta}r(x, y)\right)$, is the partition function. Details of the derivation appear in Appendix \ref{app:derivation1}. In practice, estimating the partition function $Z(x)$ remains computationally burdensome, even when we substitute the reward $r$ with its MLE estimate $r_\phi$ for the true reward $r^*$~\citep{korbak2022reinforcement, go2023aligning}, hindering practical deployment of this formulation. To address this, Eq.~\ref{eq:op_policy} can be rearranged, isolating the reward in terms of the associated optimal policy $\pi_r$, the reference policy $\pi_\text{ref}$, and the partition function $Z(\cdot)$. By applying the logarithm to both sides of Eq.~\ref{eq:op_policy} and reorganizing, we have:

\begin{equation}\label{eq:main_eq}
    r(x,y) =\beta \log \frac{\pi_r(y\mid x)}{\pi_\text{ref}(y\mid x)} + \beta \log Z(x).
\end{equation}

This formulation can equally be used for the ground-truth reward $r^*$ and the corresponding optimal policy $\pi^*$. Notably, since the Bradley-Terry model is a function of reward differences across pairs of outputs, i.e., ${p^*(y_1 \succ y_2 \mid x) = \sigma(r^*(x, y_1) - r^*(x, y_2))}$, substituting Eq.~\ref{eq:main_eq} for $r^*(x,y)$ in Eq.~\ref{eq:bradley-terry} removes the dependency on the partition function, resulting in the human preference probability being solely determined by the optimal policy $\pi^*$ and the reference policy $\pi_\text{ref}$. Therefore, under the Bradley-Terry framework, the optimal RLHF policy $\pi^*$ is characterized by the following preference probability:

\begin{equation}\label{eq:objective}
    p^*(y_1\succ y_2 \mid x)=\frac{1}{1 + \exp\left(\beta \log \frac{\pi^*(y_2\mid x)}{\pi_\text{ref}(y_2\mid x)} - \beta \log \frac{\pi^*(y_1\mid x)}{\pi_\text{ref}(y_1\mid x)}\right)}
\end{equation}

A thorough derivation is presented in Appendix~\ref{app:derivation2}. Although Eq.~\ref{eq:objective} adopts the Bradley-Terry model, analogous derivations using the broader Plackett-Luce family of models~\citep{plackett1975analysis, luce2012individual} are available in Appendix~\ref{app:plackett_luce_models}.

Given this parameterization—expressing the likelihood of human preferences in terms of the optimal policy—we can cast the maximum likelihood estimation problem over a parameterized policy $\pi_\theta$. In analogy to reward-based modeling approaches (see Eq.~\ref{eq:reward_model}), this yields the policy optimization objective:

\begin{equation}\label{eq:optimum_model}
    \mathcal{L}_\text{DPO}(\pi_{\theta}; \pi_\text{ref}) = -\mathbb{E}_{(x, y_w, y_l)\sim \mathcal{D}}\left[\log \sigma \left(\beta \log \frac{\pi_{\theta}(y_w\mid x)}{\pi_\text{ref}(y_w\mid x)} - \

Through this approach, we construct an implicit reward via an alternative parameterization such that its optimal policy corresponds directly to $\pi_\theta$. Notably, this procedure is mathematically equivalent to fitting a reparametrized Bradley-Terry model, which imparts several theoretical guarantees, such as consistency under appropriate conditions on the distribution of the preference data \cite{bong2022generalized}. Section~\ref{sec:theory} elaborates further on the theoretical aspects of DPO and contrasts them with prior works.

\textbf{How does the DPO update operate?} To gain mechanistic insight into DPO, it is helpful to examine the gradient of the loss function $\mathcal{L}_\text{DPO}$. The parameter gradient takes the following form:

\begin{multline*}\label{eq:gradient}
    \nabla_\theta \mathcal{L}_\text{DPO}(\pi_\theta;\pi_\text{ref}) = \\ -\beta\mathbb{E}_{(x, y_w, y_l) \sim \mathcal{D}} \bigg[\underbrace{\sigma(\hat{r}_\theta(x, y_l) - \hat{r}_\theta (x, y_w))}_\text{greater influence when reward ordering is incorrect}\bigg[\underbrace{\nabla_\theta\log \pi(y_w \mid x)}_\text{promote $y_w$} - \underbrace{\nabla_\theta\log\pi(y_l \mid x)}_\text{demote $y_l$}\bigg]\bigg],
\end{multline*}

where $\hat{r}_\theta(x, y) = \beta \log \frac{\pi_\theta(y \mid x)}{\pi_\text{ref}(y \mid x)}$ serves as the implicit reward defined through the relationship between the current policy $\pi_\theta$ and the reference model $\pi_\text{ref}$ (see Section~\ref{sec:theory} for additional discussion). Intuitively, the loss function gradient $\mathcal{L}_\text{DPO}$ functions to increase the probability assigned to the preferred output $y_w$, while reducing it for the non-preferred $y_l$. Crucially, each training instance is weighted by the degree to which the implicit reward model $\hat{r}_\theta$ assigns a higher value to the dispreferred completion, scaled by $\beta$, i.e., proportional to how severely the implicit reward misranks the pair, modulated by the strength of the KL regularization. Empirical evidence from our experiments indicates that this weighting is vital: omitting it (i.e., removing the weighting factor) leads to degenerate behavior in the language model, as demonstrated in Appendix Table~\ref{tab:unlikelihood_generations}.

\textbf{Overview of DPO.} The typical workflow for DPO consists of: 1) Generating completion pairs $y_1, y_2 \sim \pi_\text{ref}(\cdot \mid x)$ for each prompt $x$, assigning human preference labels to form the offline preference dataset $\mathcal{D} = \{x^{(i)}, y_w^{(i)}, y_l^{(i)}\}_{i=1}^N$; 2) training the language model $\pi_\theta$ to minimize the loss $\mathcal{L}_\text{DPO}$, given the reference model $\pi_\text{ref}$, the dataset $\mathcal{D}$, and a chosen $\beta$. 
In applied settings, it is preferable to leverage publicly available preference datasets rather than generate new completions and annotations. Since these datasets are typically sampled with $\pi^\text{SFT}$, we set $\pi_\text{ref} = \pi^\text{SFT}$ whenever this model exists. When no such $\pi^\text{SFT}$ is at hand, we fit $\pi_\text{ref}$ via maximum likelihood estimation on the preferred completions ${(x, y_w)}$; that is, ${\pi_\text{ref} = \argmax_{\pi}\mathbb{E}_{x, y_w \sim \mathcal{D}}\left[\log \pi(y_w \mid x)\right]}$. This approach seeks to reduce the distribution gap between the unknown true reference distribution and the $\pi_\text{ref}$ model employed in DPO. Additional specifics regarding implementation details and hyperparameter choices are available in Appendix~\ref{app:implementation}.

\section{Theoretical Analysis of DPO} In this section, we further elucidate the DPO methodology, supply theoretical justification, and compare the merits of DPO with actor-critic algorithms prevalent in RLHF (e.g., PPO~\cite{schulman2017proximal}).

\label{sec:theory}

\subsection{Your Language Model Is Secretly a Reward Model} DPO circumvents the need for constructing an explicit reward model and performing RL by instead utilizing a single maximum likelihood framework. Notably, the objective in Eq.~\ref{eq:main_eq} aligns with a Bradley-Terry model whose rewards are parameterized as $r^*(x, y) = \beta \log\frac{\pi^*_\theta(y \mid x)}{\pi_\text{ref}(y \mid x)}$. Consequently, training the parameterized model $\pi_{\theta}$ mirrors the procedure for reward model optimization as formulated in Eq.~\ref{eq:reward_model}, up to a reparameterization. In what follows, we develop the theory supporting this reparameterization, demonstrate that it does not limit the expressivity of the learned reward models, and establish that it permits precise recovery of the optimal policy. Our analysis begins by introducing an equivalence relation over reward functions.

\begin{definition}
We define two reward functions $r(x, y)$ and $r'(x, y)$ to be equivalent if and only if ${r(x, y)-r'(x, y) = f(x)}$ for some function $f$.     
\end{definition}

This definition clearly establishes an equivalence relation that partitions the space of reward functions into equivalence classes. The following two lemmas can then be formulated:

\begin{lemma}\label{lemma:same_prefrence} Within the Plackett-Luce framework, and notably the Bradley-Terry model, any two reward functions belonging to the same equivalence class induce identical preference distributions.
\end{lemma}

\begin{lemma}\label{lemma:same_policy}
    Within the setting of constrained RL, any pair of reward functions from the same equivalence class yield identical optimal policies.
\end{lemma}

The arguments supporting these lemmas are straightforward and are included in Appendix \ref{app:lemma1}. The result in the first lemma highlights a classical under-specification property associated with the Plackett-Luce family \cite{plackett1975analysis}. Owing to this ambiguity, it becomes necessary to enforce identifiability constraints in order to ensure meaningful MLE solutions for the reward parameterization in Eq. \ref{eq:reward_model} \cite{bong2022generalized}. According to the second lemma, the optimal policy remains invariant within each equivalence class, indicating that our learning objective is satisfied by recovering any representative reward function from the optimal class. The following theorem is proven in Appendix~\ref{app:thm1}:

\begin{theorem}\label{thm:main}
    Assuming mild technical conditions, any reward equivalence class consistent with the Plackett-Luce (or specifically, Bradley-Terry) model can be parameterized as ${r(x, y) = \beta \log \frac{\pi(y\mid x)}{\pi_\text{ref}(y\mid x)}}$ for some policy $\pi(y\mid x)$ and reference policy $\pi_\text{ref}(y \mid x)$.
\end{theorem}

\begin{sproof}
    Let $r(x, y)$ be any reward function, which gives rise to an associated optimal policy $\pi_r(y \mid x)$ as defined in Eq. \ref{eq:op_policy}. We aim to demonstrate that one can always select a reward function from the equivalence class of $r$ that admits the stated parameterization. To this end, introduce the projection operator $f$:
\begin{equation}
    f(r; \pi_\text{ref}, \beta)(x, y) = r(x, y) - \beta\log\sum_{y}\pi_\text{ref}(y\mid x)\exp\left(\frac{1}{\beta}r(x, y)\right)
\end{equation}
Here, $f$ acts by adjusting $r$ through normalization by the logarithm of the corresponding partition function, dependent only on $x$. Thus, $f(r; \pi_\text{ref}, \beta)(x, y)$ is guaranteed to lie within the equivalence class of $r(x, y)$. Substituting $r$ as prescribed by Eq.~\ref{eq:main_eq}, one obtains $f(r; \pi_\text{ref}, \beta)(x, y) = \beta \log \frac{\pi_r(y\mid x)}{\pi_\text{ref}(y\mid x)}$. Therefore, $f$ yields a member of the same equivalence class in the desired reparameterized form, indicating no loss of generality in restricting our model to this parameterization.
\end{sproof}

An alternative interpretation of Theorem~\ref{thm:main} is that it uniquely determines, within each equivalence class, the specific reward function selected by the DPO reparameterization—namely, the one that satisfies

\begin{equation}\label{eq:lag_p}
     \sum_{y}\underbrace{\pi_\text{ref}(y\mid x)\exp\left(\frac{1}{\beta}r(x, y)\right)}_{=\pi(y\mid x)\text{, using Thm.~\ref{thm:main} reparam.}} = 1,
\end{equation}

meaning that $\pi(y\mid x)$ constitutes a legitimate probability distribution, i.e., it is properly normalized and nonnegative. Examining Eq.~\ref{eq:op_policy}, we recognize that Eq.~\ref{eq:lag_p} defines the normalization, or partition function, for the optimal policy produced by the reward $r(x, y)$. A central contribution of the DPO method is its ability to introduce specific constraints on the otherwise underdetermined Plackett-Luce (particularly the Bradley-Terry) class of preference-based models, in such a way that the flexibility of reward functions remains intact, while also ensuring that the form of the optimal policy in Eq.~\ref{eq:op_policy} is explicitly computable for every prompt $x$.

\subsection{Instability of Actor-Critic Algorithms} Our framework further enables the analysis of sources of instability in standard actor-critic approaches applied in RLHF settings, such as PPO. Adopting the RLHF workflow, we concentrate on the RL fine-tuning phase described in Section~\ref{section:prelims}. Drawing parallels to the control as inference perspective \cite{levine2018reinforcement} on constrained reinforcement learning as formalized in \ref{eq:RL}, we posit a parameterized family $\pi_{\theta}(y\mid x)$ and seek to minimize $\mathbb{D}_{\text{KL}}[\pi_{\theta}(y|x) \mid \mid \pi^*(y\mid x)]$, with $\pi^*$ denoting the optimal policy determined via the reward $r_{\phi}(y, x)$ as defined in Eq.~\ref{eq:optimum_model}. Manipulating this objective, we arrive at the following maximization problem:

\begin{equation}\label{eq:AC}
    \max_{\pi_{\theta}}\mathbb{E}_{\pi_{\theta}(y\mid x)}\bigg[\underbrace{r_{\phi}(x, y) -\beta\log\sum_{y}\pi_\text{ref}(y\mid x)\exp\left(\frac{1}{\beta}r_{\phi}(x, y)\right)}_{f(r_{\phi}, \pi_\text{ref}, \beta)} - \underbrace{\beta\log\frac{\pi_{\theta}(y\mid x)}{\pi_\text{ref}(y\mid x)}}_{\text{KL}}\bigg]
\end{equation}

The objective here aligns with that targeted by previous studies \citep{ziegler2020finetuning, stiennon2022learning, bai2022training, ouyang2022training}, wherein the DPO-equivalent reward for the reward family $r_{\phi}$ is employed. Under this framework, the normalization component in $f(r_{\phi}, \pi_\text{ref}, \beta)$ can be seen as representing the soft value function associated with the reference policy $\pi_\text{ref}$. Although this normalization factor does not alter the optimal policy, omitting it may lead to elevated variance in the policy gradient estimates, thereby destabilizing the learning process. To handle this normalization, a value function can be learned, though this route is often challenging from an optimization perspective. Another strategy, utilized in earlier research, has been to normalize the reward using the outcome from a human completion as a baseline, functioning essentially as a one-sample Monte Carlo approximation to the normalization constant. By contrast, the DPO reparameterization constructs a reward function that obviates the need for explicit baselines.

\section{Experiments}
This section presents empirical assessments of DPO’s effectiveness in learning policies directly from preference data. We start by investigating, in a carefully controlled text generation context, how DPO balances the competing objectives of reward maximization and KL-divergence minimization with respect to the reference policy, particularly in comparison to popular preference optimization approaches like PPO. Subsequently, we expand our evaluation to encompass more challenging RLHF benchmarks and larger model architectures, focusing on tasks such as summarization and dialogue. The findings indicate that, with minimal hyperparameter adjustment, DPO often matches or surpasses established baselines—including RLHF with PPO and selecting the best among $N$ sampled sequences according to a trained reward function. Prior to detailing the empirical results, we outline our experimental procedures; further specifics can be found in Appendix~\ref{app:exp_details}.

\textbf{Tasks.} We evaluate our approaches on three distinct open-ended text generation tasks. Across all tasks, methods are trained to optimize a policy from a preference dataset $\mathcal{D}=\bigl\{x^{(i)}, y_w^{(i)}, y_l^{(i)}\bigr\}_{i=1}^N$. For \textbf{controlled sentiment generation}, $x$ consists of an initial segment from a movie review within the IMDb corpus \cite{maas-EtAl:2011:ACL-HLT2011}, and the policy's objective is to produce a continuation $y$ that reflects positive sentiment. To enable precise quantitative assessment, we synthetically construct preference pairs for this task by leveraging a pre-trained sentiment classifier, selecting pairs so that $p(\text{positive}\mid x,y_w)>p(\text{positive}\mid x,y_l)$. For supervised fine-tuning (SFT), GPT-2-large is trained to convergence on IMDB train split reviews (see App~\ref{app:sentiment_details} for additional methodology). In the case of \textbf{summarization}, $x$ represents a Reddit forum entry, with the policy required to summarize its main content as $y$. Our protocol follows established approaches, utilizing the Reddit TL;DR summarization dataset \citep{volske-etal-2017-tl} and preference annotations collected by \citeauthor{stiennon2022learning}. We implement SFT by fine-tuning on forum post summaries produced by humans\footnote{\url{https://huggingface.co/CarperAI/openai_summarize_tldr_sft}} using the TRLX \citep{leandro_von_werra_2023_7790115} framework for RLHF. The underlying human preference annotations were obtained by \citeauthor{stiennon2022learning} on outputs from another SFT model with a comparable training scheme. For \textbf{single-turn dialogue}, $x$ is a user's input, which could range from scientific inquiries to personal requests. The policy must generate an informative and engaging response $y$; we employ the Anthropic Helpful and Harmless dialogue dataset \citep{bai2022training}, consisting of 170,000 dialogue exchanges between a user and a virtual assistant. Each dialogue concludes with two alternative completions, generated by a large, albeit unspecified, language model and labeled according to human preference. As no SFT checkpoint is directly available for this scenario, we construct the SFT baseline by fine-tuning an open-access language model exclusively on preferred response completions.

\textbf{Evaluation.} Our empirical assessment employs two primary evaluation strategies. For controlled sentiment generation, where the true reward function (a sentiment classifier) is available, each algorithm’s performance is assessed by mapping its achieved reward versus its KL-divergence from the reference policy, forming a frontier that directly measures optimization efficacy for the constrained reward maximization objective. In contrast, since real-world tasks lack access to a known ground truth reward, we turn to a \textit{win rate} metric, comparing each method against a baseline using GPT-4 to emulate human judgment. GPT-4 provides comparative assessments of summary quality in summarization tasks and of response helpfulness in single-turn dialogue settings. The chosen baselines are the test set reference summaries in summarization and the dataset’s preferred responses in dialogue. Although prior literature supports the utility of LMs as automatic evaluators over standard metrics \citep{Chen2023ExploringTU}, we further substantiate our selection of GPT-4 for evaluation through a human study detailed in Sec.~\ref{sec:human-judgments}. Results demonstrate a high degree of concordance between GPT-4 and human raters, with human-GPT-4 agreement matching or exceeding inter-annotator consistency.

\textbf{Methods.} Besides DPO, our comparison includes several standard techniques for aligning language models with human preferences. We apply zero-shot prompting with \textbf{GPT-J} \citep{gpt-j} for summarization and 2-shot prompting with \textbf{Pythia-2.8B} \citep{biderman2023pythia} for dialogue. Furthermore, we benchmark the \textbf{SFT} model and introduce \textbf{Preferred-FT}, created via supervised fine-tuning using the preferred output $y_w$—originating from the SFT model in sentiment and summarization experiments, or from a generic language model in dialogue. The \textbf{Unlikelihood}~\citep{welleck2019neural} strategy serves as another pseudo-supervised approach, encouraging the policy to assign high probability to $y_w$ and explicitly discouraging $y_l$, with the ‘unlikelihood’ penalty modulated by a coefficient $\alpha\in[0,1]$. We also evaluate \textbf{PPO} \citep{schulman2017proximal} trained on a reward model fitted to preference data, and \textbf{PPO-GT}, which leverages access to the actual reward signal in the controlled sentiment setup. For the sentiment tasks, PPO-GT is tested in both a standard off-the-shelf configuration \cite{leandro_von_werra_2023_7790115} and a customized variant incorporating reward normalization and hyperparameter adjustments—changes also reflected in PPO when trained with learned rewards. As an additional baseline, we use \textbf{Best of $N$}, where $N$ responses are sampled from the SFT model (or Preferred-FT for dialogue) and the response with the highest predicted reward, according to a learned reward model, is chosen. While this approach can achieve strong results by decoupling reward evaluation from PPO-style optimization, it is computationally demanding, as evaluating $N$ samples per query is infeasible at test time for all but small $N$.

\subsection{How well can DPO optimize the RLHF objective?}

The standard RLHF approach involves maximizing a reward subject to a KL-divergence constraint, which serves to prevent the policy from straying too far from the reference policy. Accordingly, meaningful algorithm comparisons must jointly consider both the attained reward and the associated KL divergence; improvements in reward that come at the expense of excessive KL are not ideal. In Figure~\ref{fig:frontier-tldr-main}, we present the relationship between reward and KL for different algorithms within the sentiment analysis task. For each algorithm, we conduct several training runs, varying the policy conservativeness parameter in each run (target KL values of $\{3,6,9,12\}$ for PPO; $\beta$ values of $\{0.05,0.1,1,5\}$ and $\alpha$ values of $\{0.05,0.1,0.5,1\}$ for unlikelihood training; and differing random seeds for preferred-FT), culminating in a total of 22 experimental runs. Following every 100 training updates until the algorithms reach convergence, policies are evaluated on a set of test prompts by measuring both the mean reward under the actual reward function and the mean sequence-level KL\footnote{This refers to the cumulative sum of per-timestep KL-divergences.} with respect to the reference policy, $\text{KL}\left(\pi\mid \mid \pi_\text{ref}\right)$. Our analysis reveals that DPO generates a reward-KL frontier that is significantly more favorable than those of the other methods—achieving higher rewards for equivalent or lower levels of KL. This outcome is particularly striking for several reasons. Notably, although both DPO and PPO seek to optimize the same objective, DPO exhibits marked efficiency gains, with its reward-KL balance surpassing that of PPO across the board. Furthermore, DPO even outperforms PPO when PPO is given access to the true reward signal (PPO-GT), establishing a superior reward-KL trade-off in all cases.

\subsection{Can DPO scale to real preference datasets?}
\label{sec:dpo-real-datasets}
We proceed to assess the fine-tuning effectiveness of DPO on tasks such as summarization and single-turn dialogue. In the context of summarization, standard automated evaluation metrics like ROUGE often exhibit weak alignment with actual human preferences~\citep{stiennon2022learning}, and earlier studies have demonstrated that leveraging PPO to fine-tune language models based on human feedback produces summaries that users find more useful. To compare various techniques, we sample outputs from models on the test partition of the TL;DR summarization dataset and calculate their mean win rates against reference summaries from the test data. Sampling is performed at temperature values ranging from 0.0 to 1.0, with corresponding win rates visualized in Figure~\ref{fig:frontier-tldr-main} (right). All competing methods—DPO, PPO, and Preferred-FT—are applied as further fine-tuning to the same GPT-J SFT model\footnote{\url{https://huggingface.co/CarperAI/openai_summarize_tldr_sft}}. Our results show that at a temperature of 0.0, DPO attains an average win rate of around 61\%, surpassing PPO, which achieves approximately 57\% at its optimal temperature. Additionally, DPO secures a higher peak win rate than the best among $N$ baseline completions. It should be noted that little to no optimization was performed for DPO’s $\beta$ hyperparameter, suggesting there may be further room for improvement. Furthermore, DPO demonstrates considerably greater stability to changes in the sampling temperature relative to PPO, whose performance can deteriorate to match that of the base GPT-J at higher temperatures. Preferred-FT yields results similar to the SFT baseline, showing little improvement. A direct comparison using human assessment is presented in Section~\ref{sec:human-judgments}; in this evaluation, human raters favored DPO samples generated at temperature 0.25 over PPO samples at the same temperature in 58\% of cases.

For single-turn dialogue tasks, we conduct experiments on the portion of the Anthropic HH dataset's test split \citep{bai2022training} that features a single human-assistant exchange. For GPT-4 assessments, the preferred completions in the test set serve as references for determining win rates among competing methods. Due to the lack of a standardized SFT model for this problem, we begin with a pre-trained Pythia-2.8B, then apply Preferred-FT to finetune a reference model on selected completions, ensuring generated outputs remain in-distribution, before continuing training via DPO. Our evaluation includes the highest-scoring outcome from 128 Preferred-FT completions (having observed that the Best of $N$ metric saturates at $N=128$; see Appendix Figure~\ref{fig:best-of-n}) and a 2-shot prompt approach using the Pythia-2.8B base, with DPO matching or surpassing performance across optimal temperatures for each technique. In addition, we assess an RLHF agent trained with PPO on the Anthropic HH dataset\footnote{\url{https://huggingface.co/reciprocate/ppo_hh_pythia-6B}} and sourced from a widely recognized implementation\footnote{\url{https://github.com/CarperAI/trlx/tree/main/examples/hh}}; however, we are unable to identify prompting or sampling conditions that yield better results than the base Pythia-2.8B. Given our findings in TL;DR and the equivalence in reward functions optimized by both strategies, we take the Best of 128 as an approximate benchmark for PPO's effectiveness. In sum, DPO stands out as the sole computationally efficient approach that reliably outperforms the preferred completions within the Anthropic HH dataset, achieving results on par with or exceeding the more resource-intensive Best of 128 baseline. Furthermore, Figure~\ref{fig:dialogue-main} illustrates that DPO reaches peak effectiveness within relatively few training steps.

\subsection{Generalization to a new input distribution}

To more thoroughly assess how PPO and DPO perform when faced with a distributional shift, we analyze the policies derived from our Reddit TL;DR summarization study by applying them to a distinct dataset: news articles from the test split of CNN/DailyMail \citep{nallapati-etal-2016-abstractive}, employing the optimal sampling temperatures established in TL;DR (0 and 0.25). Table~\ref{tab:ood} details these results. We measured GPT-4’s win rate relative to the datasets' gold-standard summaries, utilizing the same GPT-4 (C) prompt as in Reddit TL;DR, except the phrase “forum post” was substituted with “news article.” On this out-of-distribution task, DPO consistently surpasses PPO by a wide margin. These findings provide preliminary evidence that DPO policies maintain generalization capabilities comparable to those of PPO policies, despite DPO not relying on the extra unlabeled Reddit TL;DR prompts that PPO leverages.

\subsection{Validating GPT-4 judgments with human judgments}
\label{sec:human-judgments}
We implement a human evaluation to ascertain the trustworthiness of GPT-4’s ratings, drawing on outputs from the TL;DR summarization experiments and employing two distinct GPT-4 prompts. The \textbf{GPT-4 (S)} (simple) prompt requests a choice of which summary most effectively conveys the key information in a post. Meanwhile, the \textbf{GPT-4 (C)} (concise) prompt extends this by also asking which summary is more concise; this addition is motivated by our observation that under the \textbf{GPT-4 (S)} prompt, GPT-4 tends to favor summaries that are lengthier and more redundant than those preferred by humans. The full text of these prompts can be found in Appendix~\ref{app:prompts}. We carry out three pairwise comparisons that include the strongest (DPO, temp. 0.25), weakest (PPO, temp. 1.0), and an intermediate (SFT, temp. 0.25) method—each pitted against the best-performing (greedily-sampled) PPO configuration—to encompass a range of output qualities. In both prompting conditions, GPT-4’s choices align with human judgments at rates comparable to inter-human agreement, indicating that GPT-4 can reliably serve as a surrogate for human evaluation (owing to a limited pool of raters, multiple human annotations are only gathered for DPO and PPO-1 matchups). Notably, the \textbf{GPT-4 (C)} prompt typically yields win rates that better reflect human preferences; thus, we adopt this prompt for reporting principal results in Section~\ref{sec:dpo-real-datasets}. Further specifics on the human evaluation, including the rater-facing web interface and the roster of human evaluators, are provided in Appendix~\ref{app:human-study}.

\section{Discussion}
Preference-based learning presents a robust and flexible approach for developing language models that are both competent and aligned with user intentions. In this work, we introduced DPO, a streamlined methodology that enables preference-based language model training without necessitating reinforcement learning. DPO circumvents the conventional method of reframing preference learning as a reinforcement learning problem; instead, it establishes a direct correspondence between language model policies and reward functions, allowing for the straightforward application of cross-entropy loss to optimize for human preferences. This method does not compromise generality and eliminates the need for complex reinforcement learning techniques. Furthermore, DPO achieves comparable or superior results to leading RLHF approaches, such as those employing PPO, even with minimal hyperparameter adjustment. Consequently, DPO significantly lowers the entry barrier for utilizing human preferences in language model training.

\textbf{Limitations \& Future Work.} The findings presented open several avenues for subsequent research. One central question concerns the generalization abilities of DPO-trained policies when faced with out-of-distribution inputs, particularly in comparison with policies trained using explicit reward signals. Preliminary evidence indicates that DPO generalizes on par with PPO-based approaches, yet systematic investigations are warranted. For instance, it remains to be explored whether self-labeling approaches leveraging DPO policies can capitalize on datasets containing unlabeled prompts. Moreover, the manifestation of reward over-optimization within the direct preference optimization paradigm requires deeper analysis, including whether the moderate drop in performance depicted in Figure~\ref{fig:dialogue-main}-right might be attributed to such over-optimization. While this study assesses models up to 6B parameters, evaluating DPO’s scalability to substantially larger, state-of-the-art models represents a promising trajectory. In terms of evaluation methodology, our observations highlight that GPT-4’s reported win rates are sensitive to the prompting protocol, motivating future research on optimizing automated assessment techniques. Beyond language models, DPO offers substantial promise for a variety of generative modeling tasks across different domains informed by preference signals.